\section{Daily Notes}

Today I tested the minimal blog structure.

\begin{enumerate}
\item Added language routes.
\item Connected a theme switcher.
\item Kept all content in LaTeX.
\end{enumerate}

Everything works without a backend.

\section{Long Read Test}

This section is intentionally long to help with scroll behavior checks. The main idea is simple: writing in LaTeX gives a predictable structure, cleaner sections, and easier navigation when the article grows. Instead of fighting visual blocks, I can focus on argument flow and narrative rhythm.

When the text gets longer, I want stable anchors, compact metadata, and an easy way to jump between sections. In practical terms that means heading-based navigation on the side, quick tag filtering, and one-click PDF export for offline reading.

\subsection{Why this format scales}

The moment a note turns into a real article, consistency becomes more important than speed. A lightweight pipeline keeps things understandable: source in \texttt{.tex}, generated HTML for web reading, and PDF for sharing. No hidden databases, no heavy editor lock-in, and no mystery transformations.

I also like that headings become semantic checkpoints. If a reader drops in from search, the page should still be easy to scan. Short paragraphs, clear section names, and practical examples make a big difference when content is technical.

\subsection{Practical writing loop}

My current loop is straightforward.

\begin{enumerate}
\item Open the article file and add one focused section.
\item Build locally and check anchors, spacing, and code examples.
\item Verify that tags are useful and not noisy.
\item Export PDF and confirm it keeps the same structure.
\end{enumerate}

This loop sounds basic, but it removes friction. Less friction means more consistent publishing.

\section{Command Line Notes}

Sometimes a tiny command snippet carries more context than a long paragraph, especially for operational writing.

\begin{quote}
\texttt{pwd}\\
\texttt{cd ~/articles}\\
\texttt{grep -R "performance" .}\\
\texttt{ls -la}
\end{quote}

The goal is not to look clever, but to keep actions reproducible. If someone can copy a command and get the same outcome, the article did its job.

\subsection{Editing principles}

I try to keep each section useful even when read in isolation. That means:

\begin{itemize}
\item state assumptions,
\item show one concrete example,
\item explain trade-offs,
\item end with one actionable conclusion.
\end{itemize}

Over time this structure creates a consistent reading experience and makes old articles easier to maintain.

\section{Closing Thoughts}

This long entry is mainly a stress test for scrolling, layout spacing, side navigation, and metadata rendering. But it also reflects the writing style I want: practical, calm, and tool-friendly.

If the reader can quickly jump by headings, inspect tags, and export a clean PDF, then the system is doing exactly what it should.
